% SHARD-ID: CA-10-LATTICE-SYMMETRY-MOTIVATION
% SHARD-TITLE: Why Lattice Symmetry Generators Are a Problem
% SHARD-SUMMARY: Opens the lattice-symmetry programme: an abstract QM system supplies time evolution, while a lattice also supplies geometry.
% SHARD-SUMMARY: Separates sourced continuum symmetry targets from intuition-driven lattice generator guesses built from local Hamiltonian densities.
% SHARD-SUMMARY: Records the source boundary after registering Weinberg and Koo-Saleur alongside Schottenloher, OAR, and lattice-fermion sources.
% SHARD-KEYWORDS: lattice symmetry, time translation, Hamiltonian density, Poincare, Lorentz, boost, source gap

\section{Why Lattice Symmetry Generators Are a Problem}
\label{sec:lattice-symmetry-motivation}

\subsection{Status, Sources, and Dependencies}

\textbf{Status: proposal shard with sourced boundary.}  The claim that a
lattice boost should look like a first moment of the Hamiltonian density is an
intuition until a local source, derivation, or test is attached.  This shard
sets the boundary for CA-10--CA-17.

Dependencies: CA-05 for symmetry as unitary/projective-unitary representation,
CA-08 for compiler output slots, and CONVENTIONS.md (h) for lattice candidate
generator signs.

% Source: references/cft/Schottenloher2008/Schottenloher2008.md:4040 -- self-adjoint A generates U(t)=exp(-itA).
% Source: references/cft/Schottenloher2008/Schottenloher2008.md:4046 -- Stone theorem for one-parameter unitary groups.
% Source: references/cft/Schottenloher2008/Schottenloher2008.md:4102 -- relativistic quantum symmetry is projective Poincare, lifted to a unitary cover.
% Source: references/cft/Schottenloher2008/Schottenloher2008.md:4113 -- translations give commuting self-adjoint P_mu.
% Source: references/cft/Schottenloher2008/Schottenloher2008.md:4117 -- P_0 is energy H and P_j are momentum components.
% Source: references/text/CFTFromLatticeFermions.txt:88 -- lattice Hamiltonian as sum of self-adjoint local terms.
% Source: references/text/OARWavelets.txt:389 -- continuous translations can be reconstructed in a scaling limit.
% Source: references/text/OARWavelets.txt:405 -- Lorentz/conformal convergence requires further work.
% Source: references/qft/SOURCES.md:9--35 -- Weinberg QFT Vol. 1 registration and targeted OCR extraction.
% Source: references/lattice-symmetry/SOURCES.md:9--32 -- original Koo-Saleur source registration and extracted line-anchor files.

\subsection{The Minimal Quantum-Mechanical Input}

An unadorned quantum system gives a Hilbert space and, once dynamics is chosen,
a self-adjoint Hamiltonian.  By Stone's theorem this is the same kind of datum as
a strongly continuous unitary representation of the additive group
\(\mathbb R\):
\[
  U(t)=\exp(-itH).
\]
Thus the one symmetry forced by the bare dynamical presentation is time
translation.  Everything beyond this requires extra structure: spatial
coordinates, locality, a metric, a lattice graph, or a continuum target.

\subsection{What the Lattice Adds}

A lattice Hamiltonian is not merely a self-adjoint operator.  It usually comes
as a decomposition
\[
  H=\sum_j h_j
\]
where each \(h_j\) is local near a site, bond, cell, or other geometric support.
The cited lattice-fermion source uses exactly this pattern: a Hamiltonian is a
sum of self-adjoint finite-support terms, together with an observable algebra
whose generators are intended as discretisations of continuum fields.

This decomposition is extra data.  Two different decompositions of the same
operator \(H\) can lead to different candidate densities and hence different
candidate symmetry generators.  Therefore the compiler target is not
\[
  \text{Hilbert space plus one operator }H,
\]
but rather
\[
  (\text{Hilbert space},\ \text{lattice geometry},\ \{h_j\},\ \text{scaling map}).
\]

\subsection{The Continuum Target}

If the intended continuum theory is Lorentz-invariant in \(d+1\) dimensions, the
target symmetry is not just \(\mathbb R\)-time translation.  It is a
projective-unitary representation of the Poincare group, lifted as necessary to
a unitary representation of a covering group.  Its translation subgroup has
self-adjoint generators
\[
  P_0=H,\qquad P_1,\ldots,P_d,
\]
where the \(P_a\) are momenta.

\textbf{Motivating question.}  If the microscopic input only supplies a lattice
Hamiltonian density \(H=\sum_j h_j\), how do we guess the lattice analogues of
\[
  P_a,\qquad M_{ab}\text{ rotations},\qquad M_{0a}\text{ boosts}?
\]

\subsection{Intuition to Be Checked}

\textbf{Intuition 1:} a rotation is a position-dependent spatial translation.
In the plane, the vector field for an infinitesimal rotation about the origin is
\[
  v(x,y)=(-y,x).
\]
Thus if we can identify local momentum densities, rotations should be spatial
first moments of those densities.

\textbf{Intuition 2:} a boost is a position-dependent time translation.  Since
time translation is generated by energy, the one-dimensional boost candidate is
a first moment of the energy density,
\[
  K_x\ \propto\ \sum_j x_j h_j.
\]

\textbf{Intuition 3:} the commutator of a boost with time translation should
produce spatial translation, up to sign and normalization.  On the lattice this
suggests that
\[
  i[H,K_x]
\]
is a candidate momentum generator.  CA-12 derives the nearest-neighbour formula
from this ansatz and records the sign convention explicitly.

\subsection{Source Boundary}

Weinberg QFT Vol. 1 is now registered with targeted OCR anchors, and the
original Koo--Saleur source is now registered from arXiv with source-TeX and PDF
text anchors.  The present shard still cites Schottenloher for continuum
symmetry definitions, OAR/wavelet sources for scaling-limit translation
recovery, and the local lattice-fermion source for rigorous free-fermion
Koo--Saleur convergence.  CA-14 cites the original Koo--Saleur paper for the
source-local lattice Virasoro prototype.
